% Part: first-order-logic
% Chapter: models-theories
% Section: theories

\documentclass[../../../include/open-logic-section]{subfiles}

\begin{document}

\olfileid{fol}{mat}{the}

\olsection{మొదటిస్థాయి సిద్ధాంతాల ఉదాహరణలు}

\begin{ex}
భాష~$\Lang L_<$లోని కఠిన రేఖీయ క్రమాల సిద్ధాంతం కింది సమితి ద్వారా
స్వీకృతీకరించబడింది:
\begin{align*}
\{\quad & \lforall[x][\lnot x < x], \\
& \lforall[x][\lforall[y][((x < y \lor y <
    x) \lor x = y)]], \\
& \lforall[x][\lforall[y][\lforall[z][((x < y
      \land y < z) \lif x < z)]]] \quad \}
\end{align*}
ఇది ఉద్దేశించిన \tetoken{నిర్మాణాలను}{!!{structure}s} సంపూర్ణంగా పట్టుకుంటుంది: ప్రతి కఠిన
రేఖీయ క్రమం ఈ స్వీకృత వ్యవస్థకు ఒక నమూనా; తిరుగుగా, $R$ అనేది సమితి
$X$పై ఒక రేఖీయ క్రమం అయితే, $\Domain M = X$, $\Assign{<}{M} = R$గా
ఉన్న నిర్మాణం~$\Struct M$ ఈ సిద్ధాంతానికి ఒక నమూనా.
\end{ex}

\begin{ex}
$\Obj 1$ (\tetoken{స్థిరసంకేతం}{!!{constant}}), $\cdot$ (ద్విస్థానిక \tetoken{ప్రమేయ సంకేతం}{!!{function}}) ఉన్న భాషలో
సమూహాల సిద్ధాంతం కింది వాటి ద్వారా స్వీకృతీకరించబడింది:
\begin{align*}
& \lforall[x][\eq[(x \cdot \Obj 1)][x]]\\
& \lforall[x][\lforall[y][\lforall[z][\eq[(x \cdot (y \cdot z))][((x
          \cdot y) \cdot z)]]]]\\
& \lforall[x][\lexists[y][\eq[(x \cdot y)][\Obj 1]]]
\end{align*}
\end{ex}

\begin{ex}
పియానో అంకగణిత సిద్ధాంతం, అంకగణిత భాష~$\Lang L_A$లోని కింది
వాక్యాల ద్వారా స్వీకృతీకరించబడింది.
\begin{align*}
& \lforall[x][\lforall[y][(\eq[x'][y'] \lif \eq[x][y])]]\\
& \lforall[x][\eq/[\Obj 0][x']]\\
& \lforall[x][\eq[(x + \Obj 0)][x]]\\
& \lforall[x][\lforall[y][\eq[(x + y')][(x + y)']]]\\
& \lforall[x][\eq[(x \times \Obj 0)][\Obj 0]]\\
& \lforall[x][\lforall[y][\eq[(x \times y')][((x \times y) + x)]]]\\
& \lforall[x][\lforall[y][(x < y \liff \lexists[z][\eq[(z' + x)][y])]]]\\
\intertext{దీనితో పాటు కింది రూపంలోని వాక్యాలన్నీ}
& (!A(\Obj 0) \land \lforall[x][(!A(x) \lif !A(x'))]) \lif \lforall[x][!A(x)]
\end{align*}
చివరి రూపంలోని వాక్యాలు అనంత సంఖ్యలో ఉంటాయి కాబట్టి ఈ స్వీకృత వ్యవస్థ
అనంతమైనది. ఈ చివరి రూపాన్ని \emph{ఆగమన పథకం} అంటారు. (నిజానికి, మనం
ఇక్కడ చూపించినదానికంటే ఆగమన పథకం కొంచెం సంక్లిష్టమైనది.)

చివరి స్వీకృతం~$<$కు ఒక \emph{స్పష్ట నిర్వచనం}.
\end{ex}

\begin{ex}
గణితపు పునాదుల్లో (అలాగే గణిత తత్వశాస్త్రంలో) శుద్ధ సమితుల సిద్ధాంతం
ముఖ్యమైన పాత్ర పోషిస్తుంది. ఒక సమితి యొక్క \tetoken{మూలకాలు}{!!{element}s} అన్నీ కూడా
శుద్ధ సమితులే అయితే ఆ సమితి శుద్ధమైనది. అందువల్ల శూన్య సమితిని
శుద్ధమైనదిగా పరిగణిస్తాం; కానీ సమితి కానిది \tetoken{ఒక మూలకంగా}{!!a{element}} ఉన్న సమితి
శుద్ధమైనది కాదు. కాబట్టి శుద్ధ సమితులు, శూన్య సమితి నుంచి మాత్రమే,
తాము సమితులు కాని వస్తువులైన ``యూర్-ఎలిమెంట్లు'' ఏవీ లేకుండా
రూపొందినవే.

కింది వాటిని శుద్ధ సమితుల సిద్ధాంతానికి ఒక స్వీకృత వ్యవస్థగా
పరిగణించవచ్చు:
\begin{align*}
& \lexists[x][\lnot \lexists[y][y \in x]]\\
& \lforall[x][\lforall[y][(\lforall[z](z \in x \liff z \in y) \lif
\eq[x][y])]]\\
& \lforall[x][\lforall[y][\lexists[z][\lforall[u][(u \in z \liff
(\eq[u][x] \lor \eq[u][y]))]]]]\\
& \lforall[x][\lexists[y][\lforall[z][(z \in y \liff \lexists[u][(z \in
        u \land u \in x)])]]]\\
\intertext{దీనితో పాటు కింది రూపంలోని వాక్యాలన్నీ} &
\lexists[x][\lforall[y][(y \in x \liff !A(y))]]
\end{align*}
మొదటి స్వీకృతం \tetoken{మూలకాలు}{!!{element}s} లేని ఒక సమితి ఉందని (అంటే,
$\emptyset$ ఉనికిలో ఉందని) చెబుతుంది; రెండవది సమితులు విస్తారతా ధర్మాన్ని
పాటిస్తాయని చెబుతుంది; మూడవది ఏ సమితులు $X$, $Y$కైనా సమితి~$\{X,Y\}$
ఉందని చెబుతుంది; నాలుగవది ఏ సమితి~$X$కైనా సమితి~$\cup X$ ఉందని
చెబుతుంది. ఇక్కడ~$\cup X$ అనేది $X$లోని మూలకాలన్నిటి సమ్మేళనం.

చివర పేర్కొన్న \tetoken{వాక్యాలను}{!!{sentence}s} సమష్టిగా \emph{అమాయక ధర్మసంగ్రహ పథకం}
అంటారు. ప్రతి~$!A(x)$కు సమితి~$\Setabs{x}{!A(x)}$ ఉందని ఇది సారాంశంగా
చెబుతుంది---తొలి చూపులో ఇది సత్యమైన, ఉపయోగకరమైన, బహుశా అవసరమైన
స్వీకృతంలా కనిపిస్తుంది. కానీ ఇది ఈ సిద్ధాంతాన్ని సంతృప్తిపరచలేనిదిగా
చేస్తుందని తేలుతుంది కాబట్టి దీనిని ``అమాయక'' పథకం అంటారు: $!A(y)$ను
$\lnot y \in y$గా తీసుకుంటే, మనకు \tetoken{వాక్యం}{!!{sentence}}
\[
\lexists[x][\lforall[y][(y \in x \liff \lnot y \in y)]]
\]
లభిస్తుంది; ఈ \tetoken{వాక్యం}{!!{sentence}} ఏ \tetoken{నిర్మాణంలోనూ}{!!{structure}} సంతృప్తి చెందదు.
\end{ex}

\begin{ex}
\emph{భాగతత్త్వం} (మీరియాలజీ)లో \emph{భాగత్వ} సంబంధం ఒక మౌలిక
సంబంధం. సమితుల సిద్ధాంతాల మాదిరిగానే, ఈ సంబంధం గురించిన వివిధ
(కొన్నిసార్లు పరస్పర విరుద్ధమైన) భావనలను స్వీకృతీకరించే భాగత్వ
సిద్ధాంతాలు ఉన్నాయి.

భాగతత్త్వ భాషలో ఒకే ద్విస్థానిక విధేయ సంకేతం~$\Obj P$ ఉంటుంది;
$\Atom{\Obj P}{x,y}$కు ``$x$, $y$లో ఒక భాగం'' అని ``అర్థం.'' ఈ
అర్థనిర్దేశాన్ని ఉద్దేశించినప్పుడు, ఈ భాష కోసం ఉన్న \tetoken{ఒక నిర్మాణాన్ని}{!!a{structure}}
\emph{భాగత్వ నిర్మాణం} అంటాం. అయితే ఒకే ద్విస్థానిక విధేయం కోసం ఉన్న
ప్రతి నిర్మాణమూ నిజంగా ఈ పేరుకు అర్హం కాదు. ``భాగత్వాన్ని'' పట్టుకునే
అవకాశం ఉండాలంటే, $\Assign{\Obj P}{M}$ కొన్ని షరతులను
సంతృప్తిపరచాలి; వాటిని భాగత్వ సిద్ధాంతానికి స్వీకృతాలుగా నిర్దేశించవచ్చు.
ఉదాహరణకు, వస్తువులపై భాగత్వం ఒక పాక్షిక క్రమం: ప్రతి వస్తువూ తనలో
తాను ఒక (\emph{క్రమేతర} భాగమైనా సరే) భాగం; వేర్వేరు రెండు వస్తువులు
ఒకదానిలో మరొకటి రెండూ భాగాలు కాలేవు; ఒక వస్తువు యొక్క భాగంలోని భాగం
కూడా ఆ వస్తువులో భాగమే. ఈ అర్థంలో ``దీనిలో భాగం'' అనేది ``దీనికి
ఉపసమితి''ని పోలి ఉంటుంది; కానీ స్వావర్తనమూ సంక్రామకమూ కాని ``దీనిలో
మూలకం'' అనే సంబంధాన్ని పోలి ఉండదని గమనించండి.
\begin{align*}
& \lforall[x][\Atom{\Obj P}{x,x}] \\
& \lforall[x][\lforall[y][((\Part{x}{y} \land \Part{y}{x})
      \lif \eq[x][y])]] \\
& \lforall[x][\lforall[y][\lforall[z][((\Part{x}{y} \land
        \Part{y}{z}) \lif \Part{x}{z})]]]\\
\intertext{అంతేకాక, ఏ రెండు వస్తువులకైనా ఒక భాగతాత్త్విక మొత్తం ఉంటుంది
  (ఆ రెండు వస్తువులూ భాగాలుగా ఉండి, ఈ విషయంలో కనిష్ఠమైన వస్తువు).}  &
\lforall[x][\lforall[y][\lexists[z][\lforall[u][(\Part{z}{u} \liff
        (\Part{x}{u} \land \Part{y}{u}))]]]]
\end{align*}
అధిభౌతిక తత్వవేత్తలు పరిశీలించే భాగత్వపు మౌలిక సూత్రాల్లో ఇవి కొన్ని
మాత్రమే. అయితే ముందుగా కొన్ని నిర్వచిత సంబంధాలను ప్రవేశపెట్టకుండా
తదుపరి సూత్రాలను సూత్రీకరించడం లేదా రాయడం త్వరగానే కష్టమవుతుంది.
ఉదాహరణకు, భాగతత్త్వంపై ఆసక్తి ఉన్న చాలామంది అధిభౌతిక తత్వవేత్తలు కింది
దానిని కూడా చెల్లుబాటు అయ్యే సూత్రంగా భావిస్తారు: వస్తువు~$x$కు క్రమ
భాగం~$y$ ఉన్నప్పుడల్లా, $y$తో ఏ భాగమూ ఉమ్మడిగా లేని భాగం~$z$ కూడా
ఆ వస్తువుకు ఉంటుంది; అంతేకాక $y$, $z$ల సంలీనం~$x$ అవుతుంది.
\end{ex}

\end{document}
